\documentclass[12pt,a4paper]{article}
\usepackage[utf8]{inputenc}
\usepackage[spanish,german]{babel}
\usepackage[T1]{fontenc}
\usepackage{geometry}
\geometry{a4paper, margin=2.5cm}
\usepackage{titlesec}
\usepackage{booktabs}
\usepackage{longtable}
\usepackage{array}
\usepackage{xcolor}
\usepackage{csquotes}

\titleformat{\section}
{\Large\bfseries}
{}
{0em}
{}[\titlerule]

\titleformat{\subsection}
{\large\bfseries}
{}
{0em}
{}

\setlength{\parindent}{0pt}
\setlength{\parskip}{0.5em}

% Farben für bessere Lesbarkeit
\definecolor{headerblue}{RGB}{0,102,204}
\definecolor{rowgray}{RGB}{245,245,245}

\begin{document}

\title{\textbf{Übersicht der erstellten Spanisch-Dialoge}}
\author{Spanisch A1/A2 -- Lernmaterialien}
\date{\today}
\maketitle

\vspace{1em}

Dieses Dokument bietet eine systematische Übersicht aller bisher erstellten Dialoge für das Spanisch-Lernen. Jeder Dialog enthält einen realistischen Gesprächstext, eine Vokabelliste, Verständnisfragen, die Konjugation eines irregulären Verbs und eine Grammatikerklärung.

\vspace{1em}

\section*{Übersichtstabelle}

\begin{longtable}{p{0.8cm}p{4.5cm}p{1.2cm}p{2.5cm}p{5.5cm}}
\toprule
\textbf{Nr.} & \textbf{Thema} & \textbf{Wörter} & \textbf{Irreguläres Verb} & \textbf{Grammatikthema} \\
\midrule
\endfirsthead

\multicolumn{5}{c}%
{{\bfseries \tablename\ \thetable{} -- Fortsetzung von vorheriger Seite}} \\
\toprule
\textbf{Nr.} & \textbf{Thema} & \textbf{Wörter} & \textbf{Irreguläres Verb} & \textbf{Grammatikthema} \\
\midrule
\endhead

\midrule \multicolumn{5}{r}{{Fortsetzung auf nächster Seite}} \\ \midrule
\endfoot

\bottomrule
\endlastfoot

01 & Saludo y presentación & 687 & SER (sein) & Das Verb \enquote{gustar} -- Gefallen ausdrücken \\
\midrule
05 & Hablar del tiempo & 756 & IR (gehen) & Das Verb \enquote{ir} -- Gehen und Zukunft (ir + a + Infinitiv) \\
\midrule
09 & Decir la hora & 723 & DECIR (sagen) & Die Uhrzeit auf Spanisch \\
\midrule
13 & Describir objetos & 734 & VER (sehen) & Adjektive zur Beschreibung von Objekten \\
\midrule
18 & En el transporte público & 812 & SALIR (ausgehen, verlassen) & Direkte Objektpronomen und Imperativ \\
\midrule
23 & Describir el lugar donde vives & 892 & TENER (haben) & Ser vs. Estar, hay/hay que \\
\midrule
33 & Hablar de salud y ejercicio & 856 & PODER (können) & Reflexive Verben, Gerundio \\
\midrule
45 & Hablar de medio ambiente & 842 & HACER (machen, tun) & Por vs. Para \\
\bottomrule
\end{longtable}

\vspace{2em}

\section*{Statistik}

\begin{itemize}
    \item \textbf{Anzahl der Dialoge:} 8
    \item \textbf{Durchschnittliche Wortanzahl:} 775 Wörter
    \item \textbf{Minimale Wortanzahl:} 687 Wörter (Dialog 01)
    \item \textbf{Maximale Wortanzahl:} 892 Wörter (Dialog 23)
    \item \textbf{Gesamtwortanzahl:} 6.202 Wörter
\end{itemize}

\vspace{2em}

\section*{Verwendete irreguläre Verben}

Die folgenden irregulären Verben wurden bereits in den Dialogen verwendet:

\begin{itemize}
    \item \textbf{SER} (sein) -- Dialog 01
    \item \textbf{IR} (gehen) -- Dialog 05
    \item \textbf{DECIR} (sagen) -- Dialog 09
    \item \textbf{VER} (sehen) -- Dialog 13
    \item \textbf{SALIR} (ausgehen, verlassen) -- Dialog 18
    \item \textbf{TENER} (haben) -- Dialog 23
    \item \textbf{PODER} (können) -- Dialog 33
    \item \textbf{HACER} (machen, tun) -- Dialog 45
\end{itemize}

\vspace{2em}

\section*{Abgedeckte Grammatikthemen}

Die Dialoge decken folgende Grammatikthemen ab:

\begin{itemize}
    \item Das Verb \enquote{gustar} und indirekte Objektpronomen
    \item Das Verb \enquote{ir} und die Futurform \enquote{ir + a + Infinitiv}
    \item Die Uhrzeit auf Spanisch
    \item Adjektive und ihre Verwendung
    \item Direkte Objektpronomen und Imperativ
    \item Ser vs. Estar und \enquote{hay/hay que}
    \item Reflexive Verben und Gerundio
    \item Por vs. Para
\end{itemize}

\vspace{2em}

\section*{Niveaustufen}

\begin{itemize}
    \item \textbf{A1 (Einfach):} Dialoge 01, 05, 09, 13
    \item \textbf{A1-A2 (Leicht-Mittel):} Dialoge 18, 23
    \item \textbf{A2 (Mittel):} Dialoge 33, 45
\end{itemize}

\vspace{2em}

\section*{Hinweise}

\begin{itemize}
    \item Alle Dialoge verwenden \textbf{südamerikanisches Spanisch}
    \item Jeder Dialog enthält:
    \begin{itemize}
        \item Einen realistischen Gesprächstext (500-1000 Wörter)
        \item Eine Vokabelliste (Spanisch-Deutsch)
        \item 20 offene Verständnisfragen
        \item Die Konjugation eines irregulären Verbs in 5 Zeiten
        \item Eine Grammatikerklärung auf Deutsch
    \end{itemize}
    \item Die Dialoge sind nach Komplexität des Vokabulars sortiert
    \item Die Verb-Konjugationen enthalten \enquote{vosotros/as} für Vollständigkeit, mit einem Hinweis, dass es in südamerikanischem Spanisch nicht verwendet wird
\end{itemize}

\end{document}
